\section{A closed shell with a moving annulus}
\label{sec:introduction}

Imagine a round disk rotating inside a stationary shell. If the disk
is a surface of revolution, its image does not change. If its height
depends on angle, the image does change, but a rotation will generally
move its attaching boundary. We arrange a third possibility: the
metric is rotational, the embedded disk is not, and a whole collar of
its boundary is planar. An intrinsic rotation followed by the opposite
ambient rotation then fixes the collar pointwise. A stationary cap
closes the disk and turns its relative rotation into a change of shape
of a closed surface.

The geometrical motion is elementary once such a disk exists. Most of
this paper is devoted to its existence, and especially to the exact
flatness required at the attaching collars.

\begin{maintheorem}\label{thm:main}
There are a smooth surface $\Sigma$ diffeomorphic to $S^2$, a smooth
Riemannian metric $g_\Sigma$, and a jointly smooth map
$\mathcal F:\Sigma\times\R\longrightarrow\R^3$ such that, writing
$F_\tau=\mathcal F(\,\cdot\,,\tau)$,
\[
 F_\tau\text{ is an embedding},\qquad
 F_\tau^*g_{\mathrm E}=g_\Sigma,\qquad F_{\tau+1}=F_\tau.
\]
At least two images $F_\tau(\Sigma)$ are not congruent under the
Euclidean isometry group. An inner disk and a complementary cap are
fixed pointwise throughout the family. The area, enclosed volume and
total oriented mean curvature are independent of $\tau$.
\end{maintheorem}

Here a closed surface is compact without boundary. A bending refers
to a family of isometric embeddings of one fixed metric; nontriviality
in this paper always means noncongruence of images, not just motion of
material points. Mean curvature is $(k_1+k_2)/2$, with the outward
orientation. These conventions matter because a relabelling of a
fixed surface is an isometric motion of its material points, but does
not change its shape.

\begin{figure}[htbp]
\centering
\begin{tikzpicture}[scale=.91,>=Stealth,font=\small]
\foreach \shift/\phase/\word in {0/0/initial,6/55/later}{
 \begin{scope}[xshift=\shift cm]
 \fill[gray!15] (0,0) circle (2.05);
 \fill[blue!6] (0,0) circle (1.65);
 \draw[gray!65,thick] (0,0) circle (2.05);
 \draw[dashed] (0,0) circle (1.65);
 \foreach \radius in {.85,1.12,1.38}{
  \draw[blue!70!black,thick,rotate=\phase]
   plot[domain=0:360,samples=65,variable=\t]
   ({(\radius+.10*cos(2*\t))*cos(\t)},
    {(\radius+.10*cos(2*\t))*sin(\t)});}
 \filldraw[fill=green!15,draw=green!45!black] (0,0) circle (.52);
 \fill[red!75!black] (1.30,1.33) circle (2.4pt);
 \node at (0,-2.36) {\word\ position};
 \end{scope}}
\draw[->,thick] (2.45,0)--(3.55,0);
\node[above] at (3,0) {phase};
\node[below,align=center] at (3,-.4) {same metric};
\node[above] at (0,2.15) {fixed inner disk};
\draw[->,green!45!black] (-.5,2.0)--(-.18,.45);
\node[above] at (6,2.15) {fixed cap marker};
\draw[->,red!75!black] (6.8,2.0)--(7.30,1.36);
\end{tikzpicture}
\caption{Relative rotation, viewed from above. The blue curves are
schematic height contours, not numerical data for the constructed
surface. The circular collars and the green inner disk stay fixed
pointwise. The cap's red marker stays fixed while the nonradial part
turns, so an ambient rotation cannot account for the entire motion.}
\label{fig:surface-motion}
\end{figure}

\subsection{Where the question sits in surface theory}

Ghomi asks whether a closed $C^2$ surface in $\R^3$ can be flexible
in Problem~1.1 of \cite{Ghomi2019}. The problem requires a continuous
family. Replacing a dimple by its reflection across a flat rim gives
two different isometric realizations, but gives no path between them.
Theorem~\ref{thm:main} provides such a path in the smooth category.

The distinction from infinitesimal bending is just as substantial.
A normal vector field supported inside a planar patch has zero first
variation of the metric. This observation already appears in
Cohn-Vossen's discussion of closed infinitesimally nonrigid surfaces
\cite[footnote~2]{CohnVossen1930}. It does not solve the nonlinear
isometry equations. Efimov's classical survey
\cite{Efimov1948Survey} and his work on deformation at zero-curvature
points \cite{Efimov1948Zero} belong to the historical setting of this
distinction; the later survey \cite[Sections~1.1--1.4]{IvanovaSabitov1991}
treats finite bendings, their infinitesimal equations and higher-order
extension. Our first deformation is also infinitesimal. The analytic
argument below is what turns it into an actual annular height.

On the rigid side, the hypotheses on curvature and boundary data
are decisive. For example, Ghomi and Spruck
\cite[Theorem~1.1]{GhomiSpruck2020} prove that two properly oriented
$C^{2,1}$ positively curved isometric immersions which agree on a
curve segment agree everywhere on their connected domain. The present
annulus has both signs of curvature, with infinitely many transitions
at either flat end. Thus neither strict positivity nor a finite
collection of nondegenerate interfaces describes its geometry.

Connelly's flexible embedded polyhedral sphere \cite{Connelly1977}
establishes flexibility in the polyhedral category. Rounding its
edges would change the metric, so it does not furnish a smooth
bending. Nash's and Kuiper's $C^1$ theorems
\cite{Nash1954,Kuiper1955} concern another regularity regime. In the
smooth problem the second fundamental form and the type of the
Darboux equation must be controlled at every scale.

Our analytic method has closer connections with smooth isometric
embedding. Smoothing compensates for derivative loss already in
Nash's argument \cite{Nash1956}. Han, Hong and Lin
\cite[Theorem~1.1]{HanHongLin2003} use a Nash--Moser--H\"ormander
scheme for nonpositive curvature under a specified degeneracy
condition. Han and Khuri treat local mixed-type problems with
finite-order changes of type \cite{HanKhuri2010,HanKhuri2013}.
In particular, the residual modification of the linearization in
\cite[pp.~650--653]{HanKhuri2010} is a precedent for keeping useful
geometric features fixed before the nonlinear equation is solved.
Here the extra issue is simultaneous control of infinitely many
interfaces and both flat ends on a single closed domain.

The three invariants in Theorem~\ref{thm:main} follow directly from
our rotation construction. They should be read with their respective
scopes: Sabitov's volume theorem \cite{Sabitov1996} is polyhedral;
Almgren and Rivin \cite{AlmgrenRivin1998} prove the smooth bending
invariance of the mean-curvature integral. Curvature-precession
motions also arise in the shell model of Hamouche, Maurini, Vidoli
and Vincenti \cite{Hamouche2017}. That mechanical comparison does
not supply the flat collars or the closed embedded surface here.

\subsection{How the construction is organized}

Let $\T=\R/(2\pi\mathbb Z)$ and choose a radial metric
\[
 g=(1+w(r)^2)\,dr^2+r^2\,d\theta^2,
 \qquad h'(r)=w(r).
\]
The radial graph of $h$ has this metric. We seek a different height
$z=h+qu$, where $q$ is a small construction parameter. The condition
that $g-dz\otimes dz$ be positive and flat is the Darboux equation.
It is enough to solve that scalar equation with a nonradial $u$ and
constant endpoint extensions. Section~\ref{sec:geometric-closure}
proves this geometric reduction completely, including injectivity
and noncongruence.

Three decisions then determine the analysis. First, a monotone flat
end kills every nonzero flat angular Jacobi field. We therefore use
an infinite sequence of rising and falling pieces, and design their
scales so that one angular mode survives while all jets still vanish
at the endpoint. Second, propagation of characteristics and matching
of endpoint solution lines have different resonances. We separate
them in a finite model before making the infinite construction.
Third, the inverse is needed for nearby nonlinear states, on the
same function space. We adjust the linearization by terms
proportional to the unsolved equations, keeping the turning circles
and characteristic returns fixed. These terms vanish at a solution.

The estimates are arranged to spend precision where these decisions
are implemented. The model spectra and elementary operator facts
are standard inputs; the matching of infinitely many intervals,
the fixed domain of the perturbed inverse, and the improvement of
endpoint decay receive the detailed proofs. In particular the
numbers $2n^3/3$, $n^{5/2}$, $12$ and $3/4$ are introduced with the
inequalities they must satisfy, rather than as unexplained choices.

Once a nonzero $q$ has been fixed, the height and the metric are fixed.
The bending parameter is a different variable:
\begin{equation}\label{eq:phase-family}
 F_\tau(r,\theta)=R_{-2\pi\tau}F(r,\theta+2\pi\tau).
\end{equation}
The rest of the proof returns to this formula. No nonlinear
integration of an infinitesimal bending in the variable $\tau$
remains to be done.
